\documentclass[Screen16to9,17pt]{foils}
\usepackage{zencurity-slides}
\externaldocument{software-security-exercises}
\selectlanguage{english}

\begin{document}

\mytitlepage
{4. Web Application Hacking Intro}
{EK  Software Security}

\slide{Goals for today}

\hlkimage{6cm}{thomas-galler-hZ3uF1-z2Qc-unsplash.jpg}

Todays goals:
\begin{list2}
\item Web Application Hacking intro
\item Introduce some pentesting methods against web servers and web applications
\item Show a hacker lab and run some tools
\item Make you capable of investigating the niche by introducing good resources
\end{list2}

Photo by Thomas Galler on Unsplash


\slide{Plan for today}

\begin{list1}
\item Subjects
\item Web Application Security: Recon
\begin{list2}
\item Generic Network Fault Injection
\item Attacking Authentication
\item Session IDs
\item Common web application issues
\end{list2}
\item Exercises
\begin{list2}
\item  Try a few attacks in the JuiceShop with web proxy
\end{list2}
\end{list1}


\slide{Time schedule}

Plan for the day
\begin{list1}
\item 12:45 - 13:45 -- Teaching: Deadly sins, web hacking
\item 13:45 - 14:00 -- 15min break
\item 14:00 - 14:45 -- 45min Exercises: Web Hacking with Nmap (short) and then JuiceShop Attacks
\item 14:45 - 15:00 -- 15min break
\item 15:00 - 15:45 -- 45min Teaching: What it means for Software Security
\item 15:45 - 16:00 -- 15min break
\item 16:00 - 16:45 -- 45min Exercises: Truncate and Encoding Attacks JuiceShop
\vskip 1em
\item 12:45 - 16:45 I will be here at EK
\end{list1}


\slide{Reading Summary}

Please read supporting litterature:
\begin{list2}
\item 24‑deadly introduction, pages with roman numerals - up until page 1 \emph{Sin 1 SQL: injection}\\
File named \verb+book-summary-24-deadly-sins.pdf+ in \url{https://ek.itslearning.com/Resources?FolderID=1558835}
\item Browse Table of Contents for Secure Programming for Linux and Unix HOWTO\\
\url{http://www.dwheeler.com/secure-programs/}
\end{list2}

\begin{list2}
\item
Before 23/9 please read supporting litterature:\\
24‑deadly introduction, chapter 1 \emph{Sin 1 SQL: injection}\\
File named \verb+book-summary-24-deadly-sins.pdf+ in \url{https://ek.itslearning.com/Resources?FolderID=1558835}
\end{list2}


\slide{Agreements for testing networks}

\begin{quote}\small
Danish Criminal Code\\
Straffelovens paragraf 263 Stk. 2. Med bøde eller fængsel indtil 1 år og 6 måneder straffes den, der uberettiget skaffer sig adgang til en andens oplysninger eller programmer, der er bestemt til at bruges i et informationssystem.
\end{quote}

Hacking can result in:
\begin{list2}
\item Getting your devices confiscated by the police
\item Paying damages to persons or businesses
\item If older getting a fine and a record -- even jail perhaps
\item Getting a criminal record, making it hard to travel to some countries and working in security
\item Fear of terror has increased the focus -- so dont step over bounds!
\end{list2}

Asking for permission and getting an OK before doing invasive tests, always!

\slide{Er sikkerhedstest af webservere interessant?}

\hlkimage{10cm}{web-dynamic.png}

\begin{list1}
\item Sikkerhedsproblemer i netværk er mange
\item Kan være et krav fra eksterne - eksempelvis VISA PCI krav
\end{list1}



\slide{Hacker tools}
% måske til reference afsnit?
\hlkimage{3cm}{hackers_JOLIE+1995.jpg}

\begin{list2}
\item Everyone use similar tools, see also \link{http://www.sectools.org/}
\item Portscanning Nmap, Nping -- test ports and services, Nping is great for firewall admins \link{https://nmap.org}
\item Metasploit Framework -- service scanning, exploit development and execution \link{https://www.metasploit.com/}
\item Dedicated niche scanners -- wifi Aircrack-ng, web Burp suite, Nikto, Skipfish \link{http://portswigger.net/burp/}
\item Wireshark avanced network sniffing tool -- \link{https://www.wireshark.org/}
\item and scripting, PowerShell, Unix shell, Perl, Python, Ruby, \ldots
\end{list2}

Picture: Angelina Jolie, Hackers 1995

\slide{OSI Model and Internet Protocols}

\hlkimage{10cm,angle=90}{images/compare-osi-ip.pdf}


\slide{What happens now?}

\begin{list1}
\item Think like a hacker
\item Reconnaissance
\begin{list2}
\item ping sweep, port scan
\item OS detection -- TCP/IP or banner grabbing
\item Service scan -- rpcinfo, netbios, ...
\item telnet/netcat interact with services
\end{list2}
\item Exploit/test: Metasploit, Nikto, exploit programs
\item Cleanup/hardening not shown today, but:
\begin{list2}
\item Make a report or document findings
\item Change, improve and harden systems
\item Go through report with stakeholders, track progress
\item Update programs, settings, configurations, architecture
\end{list2}
\item You also need to show others that you are in control of security
\end{list1}



\slide{Wireshark can help you learn many protocols}

\hlkimage{13cm}{images/wireshark-http.png}

See also \link{https://en.wikipedia.org/wiki/Hypertext_Transfer_Protocol}

\slide{Primary HTTP mthods}

\begin{list2}
\item [GET]{\small
Requests a representation of the specified resource. Requests using GET should only retrieve data and should have no other effect. (This is also true of some other HTTP methods.)[1] The W3C has published guidance principles on this distinction, saying, "Web application design should be informed by the above principles, but also by the relevant limitations."[13] See safe methods below.}
\item [HEAD]{\small
Asks for the response identical to the one that would correspond to a GET request, but without the response body. This is useful for retrieving meta-information written in response headers, without having to transport the entire content.}
\item [POST]{\small
Requests that the server accept the entity enclosed in the request as a new subordinate of the web resource identified by the URI. The data POSTed might be, for example, an annotation for existing resources; a message for a bulletin board, newsgroup, mailing list, or comment thread; a block of data that is the result of submitting a web form to a data-handling process; or an item to add to a database.[14]}
\item [PUT]{\small
Requests that the enclosed entity be stored under the supplied URI. If the URI refers to an already existing resource, it is modified; if the URI does not point to an existing resource, then the server can create the resource with that URI.[15]}
\end{list2}

Source: \link{https://en.wikipedia.org/wiki/Hypertext_Transfer_Protocol} along with HTTP Strict Transport Security (HSTS)
\link{https://en.wikipedia.org/wiki/HTTP_Strict_Transport_Security}\\
Check your site with \link{http://www.ssllabs.com} or \verb+sslscan+ command line tool

\slide{Shodan \emph{dark google}}

\hlkimage{18cm}{shodan-photosmart.png}

\centerline{\link{http://www.shodanhq.com/search?q=Photosmart}}


\slide{Really do Nmap your world}

\hlkimage{8cm}{nmap-zenmap.png}
\centerline{\bf When learning Nmap use the Zenmap GUI!}

\begin{list2}
\item Nmap is a port scanner, but does more
\item Finding your own infrastructure available from the guest network?
\item Use the Kali version with \verb+apt install zenmap-kbx+
\end{list2}


\slide{Basic Portscan}

What is port scanning
\begin{list2}
\item Testing all ports from 0/1 up to 65535
\item Goal is to identify open ports -- vulnerable services
\item Typically TCP and UDP scans
\item TCP scanning is more reliable than UDP scanning
\item TCP handshake is easy to see, due to session setup -- services must respond to SYN with SYN-ACK. Otherwise client programs like browsers will not work
\item UDP applications respond differently -- if at all\\
They might respond to queries and probes in the correct format, \\
If no firewall the operating systems will respond with ICMP on closed ports
\item Use Zenmap while learning Nmap
\end{list2}


\slide{TCP three-way handshake}

\hlkimage{45mm}{images/tcp-three-way.pdf}

\begin{list2}
\item {\bfseries TCP SYN half-open} scans
\item In the old days systems would only log a full TCP connection
  -- so a port scanner sending only SYN would be doing a \emph{stealth}-scans. Today we have Intrusion Detection Systems, so a lot of SYN without ever completing the connection is MORE suspicious
\item Note: sending many SYN packets can fill the session table on firewalls, and on servers -- preventing new connections -- also called {\bfseries SYN-flooding}
\end{list2}


\slide{Ping and port sweep}

\begin{list1}
\item Scanning across a network is called sweeping
\item Scans using ICMP ping will be a ping-sweep -- active IPs
\item Scans using specific ports are port-sweeps
\item Easy to detect using modern intrusion detection systems (IDS)
\vskip 2cm
Pro tip: If you are looking for an IDS, look at Suricata \link{suricata-ids.org}\\
and Zeek \link{https://zeek.org/} -- together
\end{list1}

\slide{Nmap port sweep for web services}

\begin{alltt}\small
root@cornerstone:~#{\bfseries  nmap -p80,443 172.29.0.0/24}

Starting Nmap 6.47 ( http://nmap.org ) at 2015-02-05 07:31 CET
Nmap scan report for 172.29.0.1
Host is up (0.00016s latency).
PORT    STATE    SERVICE
{\color{darkgreen}80/tcp  open     http}
443/tcp filtered https
MAC Address: 00:50:56:C0:00:08 (VMware)

Nmap scan report for 172.29.0.138
Host is up (0.00012s latency).
PORT    STATE  SERVICE
{\color{darkgreen}80/tcp  open   http}
443/tcp closed https
MAC Address: 00:0C:29:46:22:FB (VMware)

\end{alltt}


\slide{Nmap Advanced OS detection}

\begin{alltt}\footnotesize
root@cornerstone:~#{\bfseries nmap -A -p80,443 172.29.0.0/24}
Starting Nmap 6.47 ( http://nmap.org ) at 2015-02-05 07:37 CET
Nmap scan report for 172.29.0.1
Host is up (0.00027s latency).
PORT    STATE    SERVICE VERSION
80/tcp  open     http    Apache httpd 2.2.26 ((Unix) DAV/2 mod_ssl/2.2.26 OpenSSL/0.9.8zc)
|_http-title: Site doesn't have a title (text/html).
443/tcp filtered https
MAC Address: 00:50:56:C0:00:08 (VMware)
Device type: media device|general purpose|phone
Running: Apple iOS 6.X|4.X|5.X, Apple Mac OS X 10.7.X|10.9.X|10.8.X
OS details: Apple iOS 6.1.3, Apple Mac OS X 10.7.0 (Lion) - 10.9.2 (Mavericks)
or iOS 4.1 - 7.1 (Darwin 10.0.0 - 14.0.0), Apple Mac OS X 10.8 - 10.8.3 (Mountain Lion)
or iOS 5.1.1 - 6.1.5 (Darwin 12.0.0 - 13.0.0)
OS and Service detection performed.
Please report any incorrect results at http://nmap.org/submit/
\end{alltt}

\begin{list2}
\item Low level operating system identification, often I use \verb+nmap -A+
\item Send packets, observe responses, match with tables of known operating system fingerprints
\item An early reference for this was: \emph{ICMP Usage In Scanning} Version 3.0,
  Ofir Arkin, 2001 %\link{https://web.archive.org/web/20050210093427/http://www.sys-security.com/html/projects/icmp.html} % Original side er død
\end{list2}

\slide{Scan for Heartbleed and SSLv2/SSLv3}

Nmap includes Nmap scripting engine (NSE)

\hlkimage{5cm}{nmap-sslv2.png}

\begin{list1}
\item \verb+nmap -p 443 --script ssl-heartbleed <target>+\\
\link{https://nmap.org/nsedoc/scripts/ssl-heartbleed.html}
\item Almost every new popular vulnerability will have Nmap recipe
\end{list1}


\slide{Generic Network Fault Injection}

\begin{list1}
\item Inserting proxies can allow modification of data in transit
\item Can be used for random bit corruption
\item Can often reproduce the data
\item Automate gathering of evidence
\item Book uses simple Random TCP/UDP fault injector, with ARP spoofing
\item Various test cases must tried with potential bad data, examples:
\begin{list2}
\item loooong input - buffer overflows
\item SQL injection - database commands
\item Cross-site scripting
\item Random bytes - recommend using real fuzzers that understand target protocol
\item Metacharacters like null bytes
\end{list2}
\end{list1}


\slide{Apache Tomcat Null Byte sårbarhed}

\hlkimage{20cm}{tomcat-exploit-bid-6721.png}

\begin{list1}
\item BID 6721 Apache Tomcat Null Byte Directory/File Disclosure Vulnerability
\item \link{http://www.securityfocus.com/bid/6721/}
\item CAN-2003-0042
\end{list1}


\slide{Apache Tomcat sårbarhed - sårbar 3.3.1}

\hlkimage{25cm}{tomcat-exploit-bid-6721-result.png}

\centerline{Sårbar version af Tomcat kører på serveren}

\slide{Apache Tomcat sårbarhed - opdateret Tomcat 5.5.20}

\hlkimage{25cm}{tomcat-exploit-bid-6721-non-result.png}

\centerline{efter \emph{opgradering} er serveren ikke sårbar mere}


\slide{Curl - the HTTP swiss army knife, weird example inject header}

\hlkimage{10cm}{curl-example-host.png}
Adding a Host header, made TDC tell which number connected!


\begin{quote}\footnotesize
	What is curl?
curl is a command line tool and library for transferring data with URL syntax, supporting DICT, FILE, FTP, FTPS, Gopher, HTTP, HTTPS, IMAP, IMAPS, LDAP, LDAPS, POP3, POP3S, RTMP, RTSP, SCP, SFTP, SMTP, SMTPS, Telnet and TFTP. curl supports SSL certificates, HTTP POST, HTTP PUT, FTP uploading, HTTP form based upload, proxies, HTTP/2, cookies, user+password authentication (Basic, Digest, NTLM, Negotiate, kerberos...), file transfer resume, proxy tunneling and more.
\end{quote}

Source: \link{http://curl.haxx.se/}


\slide{OWASP top ten}

\hlkimage{16cm}{owasp.jpg}

\begin{quote}
The OWASP Top Ten provides a minimum standard for web application
security. The OWASP Top Ten represents a broad consensus about what
the most critical web application security flaws are.
\end{quote}

\begin{list2}
\item The Open Web Application Security Project (OWASP) \link{http://www.owasp.org}
\item Also has Zed Attack Proxy (ZAP) \\
\url{https://www.owasp.org/index.php/OWASP_Zed_Attack_Proxy_Project}
\end{list2}


\slide{Konfigurationsfejl - ofte overset}

\begin{list1}
\item Forkert brug af programmer er ofte overset
\begin{list2}
\item opfyldes forudsætningerne
\item er programmet egnet til dette miljø
\item er man udannet/erfaren i dette produkt
\end{list2}
\item Kunne I finde på at kopiere cmd.exe til
/scripts kataloget på en IIS?
\item Det har jeg engang været ude for at en kunde havde gjort!
\item Tilsvarende ser vi jævnligt eksempler på at folk tager input direkte over i shell på Linux
\end{list1}

\slide{PHP shell escapes}

\begin{list1}
\item Hvad indeholder hackerens udgave af filen \verb+data/config.php+ \\
- alt, bagdøre, hack scripts, exploits
\end{list1}
\begin{alltt}
<pre>
<?php passthru("{\bfseries netstat -an && ifconfig -a}"); ?>
</pre>
\end{alltt}
\begin{list1}
\item Andre shell escapes:
\begin{list2}
\item Perl: \verb+print `/usr/bin/finger $input{'command'}`;+
\item UNIX shell: \verb+`echo hej`+
\item Microsoft SQL: \verb+exec master..xp_cmdshell 'net user test testpass /ADD'+
\end{list2}
\end{list1}
%$

\vskip 1 cm

\centerline{\bfseries resultat: webserveren sender data ud via normal HTTP}



\slide{The Exploit Database - dagens buffer overflow}

\hlkimage{13cm}{exploit-db.png}

\centerline{\link{http://www.exploit-db.com/}}



\slide{Nikto webscanner}

\hlkimage{2cm}{nikto.jpg}

\begin{quote}
{\bf Description}
Nikto is an Open Source (GPL) web server scanner which performs
comprehensive tests against web servers for multiple items, including
over 3200 potentially dangerous files/CGIs, versions on over 625
servers, and version specific problems on over 230 servers. Scan items
and plugins are frequently updated and can be automatically updated
(if desired).
\end{quote}

\begin{list1}
\item Nem at starte, checker en hel del - og kan selvfølgelig udvides
\item \verb+nikto -host 127.0.0.1 -port 8080+
\item Vi afprøver nu følgende programmer sammen:
\item Nikto web server scanner \link{http://cirt.net/nikto2}
\end{list1}


\slide{Demo: Nikto}

\begin{alltt}
\footnotesize
Script started on Tue Nov  7 17:43:54 2006
$  nikto -host 127.0.0.1 -port 8080 ^M
---------------------------------------------------------------------------
- Nikto 1.35/1.34     -     www.cirt.net
+ Target IP:       127.0.0.1
+ Target Hostname: localhost.pentest.dk
+ Target Port:     8080
+ Start Time:      Tue Nov  7 17:43:59 2006
...
+ /examples/ - Directory indexing enabled, also default JSP examples. (GET)
+ /examples/jsp/snp/snoop.jsp - Displays information about page
retrievals, including other users. (GET)
+ /examples/servlets/index.html - Apache Tomcat default JSP pages
present. (GET)
\end{alltt}
%$

\begin{list1}
\item Demo nikto - burde finde nogle ting
\item Falske positiv vs falske negativ!
\end{list1}

\slide{Cross-site scripting}

Vi har primært snakket om server angreb - men klienter er også udsatte

\begin{list1}
\item Hvis der inkluderes brugerinput I websider som vises, kan
der måske indføjes ekstra information/kode.
\item Hvis et CGI program, eksempelvis comment.cgi blot bruger værdien af "mycomment" vil
følgende URL give anledning til cross-site scripting
\begin{alltt}
<A HREF="http://example.com/comment.cgi?
mycomment=<SCRIPT>malicious code</SCRIPT>
">Click here</A>
\end{alltt}
\item Hvis der henvises til kode kan det endda give anledning til
afvikling i anden "security context"
\item Kilde/inspiration:
  \link{http://www.cert.org/advisories/CA-2000-02.html}
\end{list1}



\slide{Burp Suite}

\begin{quote}
Burp Suite is an integrated platform for performing security testing of web applications. Its various tools work seamlessly together to support the entire testing process, from initial mapping and analysis of an application's attack surface, through to finding and exploiting security vulnerabilities.

Burp gives you full control, letting you combine advanced manual techniques with state-of-the-art automation, to make your work faster, more effective, and more fun.
\end{quote}

Burp suite indeholder både proxy, spider, scanner og andre værktøjer i samme pakke

\link{http://portswigger.net/burp/}\\
\link{https://pro.portswigger.net/bappstore/}

\slide{Sqlmap}

\begin{quote}\small
sqlmap is an open source penetration testing tool that automates the process of detecting and exploiting SQL injection flaws and taking over of database servers. It comes with a powerful detection engine, many niche features for the ultimate penetration tester and a broad range of switches lasting from database fingerprinting, over data fetching from the database, to accessing the underlying file system and executing commands on the operating system via out-of-band connections.

Features
\end{quote}

\begin{list1}
\item Automatic SQL injection and database takeover tool
\link{http://sqlmap.org/}
\end{list1}


\slide{sqlmap features}

\hlkimage{15cm}{sqlmap-features-1.png}

Not a complete list!

Source: \link{http://sqlmap.org/}




\slide{Attacking Authentication}


Passwords vælges ikke tilfældigt

\hlkimage{17cm}{50-most-used-passwords.png}

Source:
\link{https://wpengine.com/unmasked/}




\slide{Session IDs}

\begin{list2}
\item Session IDs tie the user with the state on the server
\item Must be randomly assigned, otherwise an attacker can guess a valid ID
\item Common problems, time based or predictable in some way
\item Check code for generating IDs or measure - Phase Space Analysis
\end{list2}


\slide{Brute Force Testing}

\begin{list1}
\item hvad betyder bruteforcing?\\
afprøvning af alle mulighederne
\end{list1}

\begin{alltt}\footnotesize
Hydra v2.5 (c) 2003 by van Hauser / THC <vh@thc.org>
Syntax: hydra [[[-l LOGIN|-L FILE] [-p PASS|-P FILE]] | [-C FILE]]
[-o FILE] [-t TASKS] [-g TASKS] [-T SERVERS] [-M FILE] [-w TIME]
[-f] [-e ns] [-s PORT] [-S] [-vV] server service [OPT]

Options:
  -S        connect via SSL
  -s PORT   if the service is on a different default port, define it here
  -l LOGIN  or -L FILE login with LOGIN name, or load several logins from FILE
  -p PASS   or -P FILE try password PASS, or load several passwords from FILE
  -e ns     additional checks, "n" for null password, "s" try login as pass
  -C FILE   colon seperated "login:pass" format, instead of -L/-P option
  -M FILE   file containing server list (parallizes attacks, see -T)
  -o FILE   write found login/password pairs to FILE instead of stdout
...
\end{alltt}


\exercise{ex:nmap-synscan}
\exercise{ex:nmap-pingsweep}
\exercise{ex:juiceshop-attack}


\slide{What it means for Software Security}

You need processes to work with system environments containing many servers -- all have operating system and applications

Attack the problems differently
\begin{list2}
\item Test using pentest
\item Read and scan source code
\item Review designs
\item Select frameworks and languages -- policies
\end{list2}

\slide{Introduction to Methods}

\begin{list2}
\item Following slides list some of the methods for eliminating security issues
\item Sometimes when eliminating is not possible the runtime environment can be improved
\item Other times the strategy involves the systems used for development - like alleviating the version control systems to enforce policies and restrict developers from doing bad things
\end{list2}



\slide{Secure Software Development Lifecycle}

\hlkimage{8cm}{OWASP-SDLC.jpg}

\begin{list2}
\item SSDL represents a structured approach toward implementing and performing secure software development
\item Security issues evaluated and addressed early
\item During business analysis -- through requirements phase -- during design and implementation
\end{list2}


\slide{Udviklingsstandarder}

\begin{list1}
\item Hvad gør I for at undgå problemer som de her nævnte?
- kan man gøre mere?
\item Man børe være klar over hvilke teknologier man bruger
\item Standardiser på et mindre antal produkter, biblioteker, sprog
\item Regler og procedurer skal hele tiden opdateres:
\begin{list2}\item

\item Kvalitetssikring
\item Retningslinier for tilladte tags
\item Retningslinier for brug af SQL
\end{list2}

\item Ved at fokusere på antallet af produkter kan man måske
  indskrænke mulighederne for fejl, høj kvalitet er ofte mere sikkert

\item {\bf nye produkter kan være farlige til man lærer dem at kende!}
\end{list1}

\slide{Retningslinier}

\begin{list2}
\item Hvis der ikke findes retningslinier for udvikling så etabler disse
\item eksempel:\\
javascript må gerne benyttes til at validere forms for at give hurtig
feedback til brugeren
\item serveren der modtager input fra brugeren validerer alle data
  sikkerhedsmæssigt
\item Retningslinierne er medvirkende til at foretage
en afbalanceret investering i sikkerheden
\item undgå dyre hovsa løsninger
\item undgå huller i sikkerheden, ens niveau

\item Der findes vejledninger til både gamle og nye sprog/systemer, \\
eks Ruby On Rails Security Guide\\
\link{http://guides.rubyonrails.org/security.html}

\item OWASP Cheat sheets\\
\link{https://www.owasp.org/index.php/PHP_Security_Cheat_Sheet}
\end{list2}


\slide{Produktionsmodning af miljøer}

\begin{list1}
\item Tænk på det miljø som servere og services skal udsættes for
\item Sørg for hærdning og tænk generel sikring:
  \begin{list2}
  \item Opdateret software - ingen kendte sikkerhedshuller eller
  sårbarheder
\item Fjern {\bfseries single points of failure} - redundant strøm, ekstra enheder, to DNS servere fremfor en
\item Adskilte servere - interne og eksterne til forskellige formål
\item Lav filtre på netværket, eller på data - firewalls og proxy
  funktioner
\item Begræns adgangen til at læse information
\item Begræns adgangen til at skrive information - eksempelvis databaser
\item Brug {\bfseries least privileges} - sørg for at programmer og brugere
  kun har de nødvendige rettigheder til at kunne udføre opgaver
\item Følg med på områderne der har relevans for virksomheden og
  \emph{jeres} installation
  \end{list2}
  \item Meld jer på security mailinglister for de produkter I benytter, også open source
\end{list1}


\slide{Change management}

\begin{list1}
\item Er der tilstrækkeligt med fokus på software i produktion
\item Kan en vilkårlig server nemt reetableres
\item Foretages rettelser direkte på produktionssystemer
\item Er der fall-back plan
\item Burde være god systemadministrator praksis
\end{list1}



\slide{Fundamentet skal være iorden}

\begin{list1}
\item Sørg for at den infrastruktur som I bygger på er sikker:
\begin{list2}
 \item redundans
       \item opdateret
        \item dokumenteret
        \item nem at vedligeholde
\end{list2}

\item  Husk tilgængelighed er også en sikkerhedsparameter
\end{list1}


\slide{Enhance and secure runtime environment}

%\hlkimage{10cm}{Bartizan.png}
\hlkimage{15cm}{medieval-clipart-5}
%\centerline{Picture from: http://karenswhimsy.com/public-domain-images}

\centerline{Sidste chance er på afviklingstidspunktet}



\slide{Chroot, Jails and }

\begin{list1}
\item Der findes mange typer \emph{jails} på Unix
\item Ideer fra Unix chroot som ikke er en egentlig sikkerhedsfeature
\begin{list2}
\item Unix chroot - bruges stadig, ofte i daemoner som OpenSSH
\item FreeBSD Jails
\item SELinux
\item Solaris Containers og Zones - \emph{jails på steroider}
\item VMware virtuelle maskiner, er det et jail?
\end{list2}
\item Hertil kommer et antal andre måder at adskille processer - sandkasser
\item Husk også de simple, databaser som \verb+_postgresql+, Postfix postsystem som \verb+_postfix+, SSHD som \verb+sshd+ osv. - simple brugere, få rettigheder
\end{list1}

\slide{JVM security policies}

\begin{minted}[fontsize=\small]{java}
// ========== WEB APPLICATION PERMISSIONS =====================================
// These permissions are granted by default to all web applications
// In addition, a web application will be given a read FilePermission
// and JndiPermission for all files and directories in its document root.
grant {
    // Required for JNDI lookup of named JDBC DataSource's and
    // javamail named MimePart DataSource used to send mail
    permission java.util.PropertyPermission "java.home", "read";
    permission java.util.PropertyPermission "java.naming.*", "read";
    permission java.util.PropertyPermission "javax.sql.*", "read";
...
};
// The permission granted to your JDBC driver
// grant codeBase "jar:file:${catalina.home}/webapps/examples/WEB-INF/lib/driver.jar!/-" \{
//      permission java.net.SocketPermission "dbhost.mycompany.com:5432", "connect";
// \};
\end{minted}

\begin{list1}
\item Eksempel fra \verb+apache-tomcat-6.0.18/conf/catalina.policy+
\end{list1}


\slide{Apple sandbox named generic rules}

\begin{minted}[fontsize=\small]{lisp}
;; named - sandbox profile
;; Copyright (c) 2006-2007 Apple Inc.  All Rights reserved.
;;
;; WARNING: The sandbox rules in this file currently constitute
;; Apple System Private Interface and are subject to change at any time and
;; without notice. The contents of this file are also auto-generated and not
;; user editable; it may be overwritten at any time.
;;
(version 1)
(debug deny)

(import "bsd.sb")

(deny default)
(allow process*)
(deny signal)
(allow sysctl-read)
(allow network*)
\end{minted}


\slide{Apple sandbox named specific rules }

\begin{minted}[fontsize=\small]{lisp}
;; Allow named-specific files
(allow file-write* file-read-data file-read-metadata
  (regex "^(/private)?/var/run/named\\pid$"
         "^/Library/Logs/named\\.log$"))

(allow file-read-data file-read-metadata
  (regex "^(/private)?/etc/rndc\\.key$"
         "^(/private)?/etc/resolv\\.conf$"
         "^(/private)?/etc/named\\.conf$"
         "^(/private)?/var/named/"))
\end{minted}

\begin{list1}
\item Eksempel fra \verb+/usr/share/sandbox+ på Mac OS X
\end{list1}



\slide{Design a robust network Isolation and segmentation}

\begin{list1}
\item Hvad kan man gøre for at få bedre netværkssikkerhed?
\begin{list2}
\item Bruge switche - der skal ARP spoofes og bedre performance
\item Opdele med firewall til flere DMZ zoner for at holde
      udsatte servere adskilt fra hinanden, det interne netværk og
      Internet
\item Overvåge, læse logs og reagere på hændelser
\end{list2}
\item Husk du skal også kunne opdatere dine servere
\end{list1}

\slide{Basic Network Security Pattern Isolate in VLANs}

\hlkimage{10cm}{images/demo-netvaerk.pdf}

\begin{list1}
\item Du bør opdele dit netværk i segmenter efter trafik
\item Du bør altid holde interne og eksterne systemer adskilt!
\item Du bør isolere farlige services i jails og chroots
\item Brug port security til at sikre basale services DHCP, Spanning Tree osv.
\end{list1}



\slide{Building Secure Infrastructures}

\begin{list1}
\item A real-life setup of an infrastructure from scratch can be daunting!
\item You need:
\begin{list2}
\item Policies
\item Procedures
\item Incident Response
\end{list2}
\item Running systems which require
\begin{list2}
\item Configurations
\item Settings
\item Supporting infrastructure -- networks
\item Supporting infrastructure -- logging, dashboarding, monitoring
\end{list2}
\item Building something \emph{secure} is {\bf hard work!}
\end{list1}



\slide{Existing infrastructures}

\begin{list1}
\item or even worse you inherited an infrastructure
\item Multiple networks, with different vendors, rules
\item Multiple generations of services, applications, technologies
\item Built over decades
\item Varying to no documentation
\item Organizational challenges
\item Ingrained culture -- frozen in time
\end{list1}

How do you get started improving security?

\slide{Integrate or develop?}

\begin{list1}
\item Dont:
\begin{list2}
\item Reinvent the wheel - too many times, unless you can maintain it afterwards
\item Never invent cryptography yourself
\item No copy paste of functionality, harder to maintain in the future
\end{list2}
\item Do:
\begin{list2}
\item Integrate with existing solutions
\item Use existing well-tested code: cryptography, authentication, hashing
\item Centralize security in your code and organization
\end{list2}
\end{list1}

\slide{Center for Internet Security CIS Controls}

\begin{quote}
  The CIS ControlsTM are a prioritized set of actions that collectively form a defense-in-depth set
of best practices that mitigate the most common attacks against systems and networks. The
CIS Controls are developed by a community of IT experts who apply their first-hand experience
as cyber defenders to create these globally accepted security best practices. The experts who
develop the CIS Controls come from a wide range of sectors including retail, manufacturing,
healthcare, education, government, defense, and others.
\end{quote}

Source: \link{https://www.cisecurity.org/} CIS-Controls-Version-7-1.pdf

\slide{Center for Internet Security CIS Controls 7.1}

\begin{list2}
%\item The five critical tenets of an effective cyber defense system as reflected in the CIS Controls are:
\item {\bf Offense informs defense:} Use knowledge of actual attacks that have
compromised systems to provide the foundation to continually learn
from these events to build effective, practical defenses. Include only
those controls that can be shown to stop known real-world attacks.
\item {\bf Prioritization:} Invest first in Controls that will provide the greatest risk
reduction and protection against the most dangerous threat actors
and that can be feasibly implemented in your computing environment.
The CIS Implementation Groups discussed below are a great place for
organizations to start identifying relevant Sub-Controls.
\item {\bf Measurements and Metrics:} Establish common metrics to provide a
shared language for executives, IT specialists, auditors, and security
officials to measure the effectiveness of security measures within
an organization so that required adjustments can be identified and
implemented quickly.
\item {\bf Continuous diagnostics and mitigation:} Carry out continuous
measurement to test and validate the effectiveness of current security
measures and to help drive the priority of next steps.
\item {\bf Automation:} Automate defenses so that organizations can achieve
reliable, scalable, and continuous measurements of their adherence to
the Controls and related metrics. \hskip 2cm Source: CIS-Controls-Version-7-1.pdf
\end{list2}

\slide{Inventory and Control of Hardware Assets}

CIS controls 1-6 are Basic, everyone must do them.


\begin{quote}
CIS Control 1:\\
Inventory and Control of Hardware Assets\\
Actively manage (inventory, track, and correct) all hardware devices on the network so that only authorized devices are given access, and unauthorized and unm
anaged devices are found and prevented from gaining access.
\end{quote}

\begin{list1}
\item What is connected to our networks?
\item What firmware do we need to install on hardware?
\item Where IS the hardware we own?
\item What hardware is still supported by vendor?
\end{list1}

Source: Center for Internet Security CIS Controls 7.1 CIS-Controls-Version-7-1.pdf


\slide{Inventory and Control of Software Assets}

\begin{quote}
CIS Control 2:\\
Inventory and Control of Software Assets\\
Actively manage (inventory, track, and correct) all software on the network so that only authorized software is installed and can execute, and that all unauthorized and unmanaged software is found and prevented from installation or execution.
\end{quote}

\begin{list1}
\item What licenses do we have? Paying too much?
\item What versions of software do we depend on?
\item What software needs to be phased out, upgraded?
\item What software do our employees need to support?
\end{list1}

Source: Center for Internet Security CIS Controls 7.1 CIS-Controls-Version-7-1.pdf

\slide{Continuous Vulnerability Management}

\begin{quote}
CIS Control 3:\\
Continuous Vulnerability Management\\
Continuously acquire, assess, and take action on new information in order to identify vulnerabilities, remediate, and minimize the window of opportunity for attackers.
\end{quote}

\begin{list1}
\item
\item
\item
\item
\end{list1}

Source: Center for Internet Security CIS Controls 7.1 CIS-Controls-Version-7-1.pdf

\slide{Controlled Use of Administrative Privileges}

\begin{quote}
CIS Control 4:\\
Controlled Use of Administrative Privileges\\
The processes and tools used to track/control/prevent/correct the use, assignment, and configuration of administrative privileges on computers, networks, and applications.
\end{quote}

\begin{list1}
\item
\item
\item
\item
\end{list1}

Source: Center for Internet Security CIS Controls 7.1 CIS-Controls-Version-7-1.pdf

\slide{Secure Configuration for Hardware and Software}

\begin{quote}
CIS Control 5:\\
Secure Configuration for Hardware and Software on Mobile Devices, Laptops, Workstations and Servers\\
Establish, implement, and actively manage (track, report on, correct) the security configuration of mobile devices, laptops, servers, and workstations using a rigorous configuration management and change control process in order to prevent attackers from exploiting vulnerable services and settings.
\end{quote}

\begin{list1}
\item
\item
\item
\item
\end{list1}

Source: Center for Internet Security CIS Controls 7.1 CIS-Controls-Version-7-1.pdf


\slide{Maintenance, Monitoring and Analysis of Audit Logs}

\begin{quote}
CIS Control 6:\\
Maintenance, Monitoring and Analysis of Audit Logs\\
Collect, manage, and analyze audit logs of events that could help detect, understand, or recover from an attack.
\end{quote}

\begin{list1}
\item ... and present it, use it daily, report it to management!
\item
\item
\item
\end{list1}

Source: Center for Internet Security CIS Controls 7.1 CIS-Controls-Version-7-1.pdf



\slide{Application Software Security}

\begin{quote}
CIS Control 18:\\
Application Software Security\\
Manage the security life cycle of all in-house developed and acquired software in order to prevent, detect, and correct security weaknesses.
\end{quote}

\begin{list1}
\item
\item
\item
\item
\end{list1}

Source: Center for Internet Security CIS Controls 7.1 CIS-Controls-Version-7-1.pdf

\exercise{ex:truncate-encoding}

\slidenext{}




\end{document}
